This chapter surveys the work our project builds on. A recurring theme is the need to
keep three different threat classes apart: \emph{physical side-channel analysis} (the
subject of this project), \emph{mathematical cryptanalysis} (which attacks the public
transcript, with no physical leakage), and \emph{fault attacks} (which actively perturb
a computation). Conflating them would misrepresent what is actually known about MQOM.

\section{MQOM and the MPCitH / TCitH Paradigm}

MQOM~\cite{mqomspec} compiles a zero-knowledge proof of knowledge of a solution to an
MQ system into a signature via the Fiat--Shamir transform. The secret witness $x$ is
additively shared; a GGM seed tree expands one root seed into $N$ leaf seeds, each
producing one share, and an all-but-one opening reveals every share except one. A
\emph{correlated} half-tree optimisation ties the first $\lambda$ bits of $x$ into the
tree, and a correction value $\Delta x = x \oplus \mathrm{acc\_x}$ (where
$\mathrm{acc\_x}$ is the sum of the shares) closes the gap between the shares and the
true secret. These structural facts, taken from the specification, determine where the
secret meets public data and hence where CPA is even possible.

\section{Side-Channel Precedents on MPCitH and MQ Signatures}

The closest physical precedent is the single-trace attack of Godard et
al.~\cite{godard2024} on the MPCitH framework, which mounts a Soft-Analytical
Side-Channel Attack (belief propagation over a factor graph) on the finite-field
multiplication inside a TCitH-family polynomial evaluation. Importantly, the authors
note that when a public multiplier row consists only of zeros and ones, even perfect
knowledge of the Hamming weight cannot isolate a single key candidate --- an early
statement of the \emph{identifiability} obstacle that is central to our work.

For multivariate schemes with a public-matrix-times-secret structure, Park et
al.~\cite{park2018} demonstrated CPA with Hamming-weight and Hamming-distance models on
the matrix-vector product in Rainbow and UOV. Their setting differs from MQOM in a
decisive way: in UOV/Rainbow the \emph{known} operand varies across signatures, giving
classical vertical CPA, whereas in MQOM's analogous multiplication both the public
matrix and the secret are fixed across signatures.

\section{The Identifiability Problem on Linear Field Maps}

A recent non-profiled CPA study on MAYO~\cite{lampropoulou2026} gives the sharpest
published statement of the difficulty that linear (as opposed to S-box-like) targets
pose. The multiplicative-identity hypothesis aliases onto input-load leakage, so a
global-maximum CPA systematically returns the wrong value; the authors remove the alias
by partial correlation (regressing the Hamming weight of the input out of the traces
before correlating). They further observe that increasing the number of traces merely
reduces the variance of an estimator that is converging on the wrong maximum --- more
data does not cure a model mismatch on a linear map. This directly informs our
treatment of MQOM's (linear) correction-value and matrix-multiply targets.

\section{CPA Methodology and Leakage Models}

The correlation distinguisher with a Hamming-weight or Hamming-distance leakage model
was introduced by Brier et al.~\cite{brier2004cpa}. The stochastic (linear-regression)
leakage model of Schindler et al.~\cite{schindler2005stochastic} estimates per-bit
leakage coefficients rather than assuming a perfect Hamming-weight model, and is the
natural tool for weighted-leakage and model-mismatch studies. Horizontal correlation
analysis~\cite{clavier2010horizontal} exploits many operations within a single trace,
which is the relevant mode when the secret meets varying public data \emph{within} one
signature rather than across signatures. For completeness we also note template
attacks~\cite{chari2002templates}, SASCA~\cite{veyrat2014sasca}, and the
information-theoretic evaluation framework of Standaert et al.~\cite{standaert2009}.

\section{Mathematical Cryptanalysis, Fault Attacks, and Masking of MQOM}

Several recent results concern MQOM but are \emph{not} side-channel attacks and must be
kept categorically separate. Delgado's passive key-recovery analysis~\cite{delgado2026}
recovers signing-key information from the public transcript by exploiting a saltless
root-expansion in the v2 lineage --- a protocol-level mathematical attack with no
physical leakage. (A companion single-signature structural claim was withdrawn by its
author and is therefore not relied upon here.) Feneuil and
Rivain~\cite{feneuil2026ggm}, and Kosuge and Xagawa~\cite{kosuge2026formal}, provide
formal security analyses. On the physical side but via a different mechanism, Sarde et
al.~\cite{sarde2025fault} mount a \emph{fault} attack on MQOM's MQ evaluation, and
Jendral and Dubrova~\cite{jendral2025voleith} analyse the VOLEitH scheme FAEST (adjacent
to, but distinct from, MQOM's TCitH line). Finally, a high-order masking of MQOM~v2.1
signing~\cite{mqom2026masking} independently confirms the correction/multiplication
operations as the sensitive, secret-dependent targets, which supports our choice of
attack points.

\section{Gap Addressed by This Project}

No end-to-end power-analysis key recovery on the MQOM~v2.1 TCitH reference has been
published. The precedents above establish that field multiplication and
public-matrix-times-secret products are plausible CPA targets, and that linear targets
raise genuine identifiability problems. This project addresses that gap in a
computer-only setting: it builds independent models of MQOM~v2.1's secret-dependent
intermediates, formulates CPA variants and leakage hypotheses, and determines by
simulation and analysis whether and how the full secret can be recovered.
